\begin{figure}[htbp]
\centering
\begin{tikzpicture}

% ---- campos del mensaje ----
\node[campo, minimum width=17mm, fill=azuloscuro!8] (acct) {1234};
\node[campo, minimum width=11mm, right=0pt of acct, fill=grisclaro] (mt) {18};
\node[campo, minimum width=8mm,  right=0pt of mt,   fill=ambar!15]  (q)   {1};
\node[campo, minimum width=14mm, right=0pt of q,    fill=verdeacento!12] (xyz) {381};
\node[campo, minimum width=11mm, right=0pt of xyz,  fill=grisclaro] (gg)  {00};
\node[campo, minimum width=14mm, right=0pt of gg,   fill=verdeacento!12] (ccc) {007};
\node[campo, minimum width=8mm,  right=0pt of ccc,  fill=grisclaro] (s)   {B};

% ---- nombres de campo encima ----
\foreach \n/\l in {acct/ACCT, mt/MT, q/Q, xyz/XYZ, gg/GG, ccc/CCC, s/S}
  \node[font=\scriptsize\bfseries\color{azuloscuro}, above=1pt of \n] {\l};

% ---- descripciones debajo, escalonadas ----
\node[nota, below=17mm of s.south, xshift=6mm, anchor=west]   (ds)   {checksum};
\node[nota, below=24mm of ccc.south, xshift=6mm, anchor=west] (dccc) {zona 007: \textbf{el dispositivo}};
\node[nota, below=31mm of gg.south, xshift=6mm, anchor=west]  (dgg)  {grupo o partición};
\node[nota, below=38mm of xyz.south, xshift=6mm, anchor=west] (dxyz) {código de evento 381: \textbf{pérdida de supervisión RF}};
\node[nota, below=45mm of q.south, xshift=6mm, anchor=west]    (dq)   {calificador 1: \textbf{nuevo evento} (3 = restauración)};
\node[nota, below=52mm of mt.south, xshift=6mm, anchor=west]  (dmt)  {tipo de mensaje, siempre 18};
\node[nota, below=59mm of acct.south, xshift=6mm, anchor=west](dacct){cuenta 1234: \textbf{el sitio}};

\foreach \c/\d in {s/ds, ccc/dccc, gg/dgg, xyz/dxyz, q/dq, mt/dmt, acct/dacct} {
  \draw[flecha, dashed, line width=0.4pt] (\c.south) |- (\d.west);
}

% ---- lectura ----
\node[cajaacento, below=6mm of dacct.south west, anchor=north west, text width=12.5cm,
      fill=verdeacento!6, font=\footnotesize\itshape]
  {«En la cuenta 1234 se produjo un nuevo evento de pérdida de supervisión RF
    en el dispositivo 007.»};

\end{tikzpicture}
\caption{Estructura del mensaje Contact ID: 16 dígitos distribuidos en siete campos.}
\label{fig:mensaje}
\end{figure}
